\section{Vibe-Artifact Taxonomy and Stable Rule IDs}

Vibe coding is not a mood. It is a fault class: shipping code whose correctness depends on plausibility, confidence, or reviewer intuition instead of machine-checkable proof. A vibe artifact is the repository residue left behind by that workflow. The risk-priority number below is \textbf{RPN = impact $\times$ likelihood $\times$ detection difficulty}, with each factor scored 1--5. It is a Jankurai policy weight, not a universal incident statistic. Future releases should calibrate it against repair success, escaped defects, security regressions, and review time.

Public evidence supports treating vibe coding as a high-risk production pattern rather than a harmless workflow style. Veracode reports that AI-generated samples introduced risky security flaws in 45\% of tests \cite{veracodeGenaiCodeSecurity2025}. GitGuardian reports accelerating hardcoded-secret exposure in public GitHub activity and AI-service leak growth \cite{gitguardianSecretsSprawl2026,gitguardianSecretsSprawl2026Blog}. Stack Overflow's 2025 survey reports that developers' largest AI frustration is output that is almost right but not quite, with time-consuming AI-code debugging close behind \cite{stackOverflowAISurvey2025}. OWASP's LLM Top 10 names prompt injection, sensitive information disclosure, supply-chain risk, improper output handling, and excessive agency as first-order AI application risks \cite{owaspLlmTopTen2025}. SetupBench and ContextBench add the operational failure modes: agents still struggle with environment bootstrap and useful context retrieval \cite{setupBench2025,contextBench2026}.

The v0.8.0 coverage method is mechanical: parse source rows, normalize issue families, map each row to one HLT rule, assign TLR/RPN, mark detector maturity, and emit registry/report artifacts. Every parsed row in the vibe-coding corpus maps to exactly one canonical issue in \texttt{agent/vibe-coverage.toml}. Detector maturity uses four labels: deterministic detectors can be reproduced from repository state; evidence-routed controls require project proof artifacts; governance-backed controls require accepted policy or review receipts; research rows are not yet detector-backed. Appendix~A renders the generated summary: green means detector-backed Jankurai control plus proof/report evidence, yellow means a reviewed control family exists but still needs issue-specific fixture, CI, or governance evidence. The v0.8.0 registry has 0 unmapped rows.

\begin{table*}[t]
\caption{Ranked vibe-artifact taxonomy. RPN is a Jankurai policy priority, not an empirical frequency claim.}
\label{tab:fault-taxonomy}
\centering
\scriptsize
\begin{tabularx}{\textwidth}{L{0.12\textwidth} R{0.04\textwidth} L{0.16\textwidth} L{0.15\textwidth} Y L{0.13\textwidth}}
\toprule
\JKTableHeader
\textbf{TLR} & \textbf{RPN} & \textbf{Fault} & \textbf{Rule/lane} & \textbf{Required control} & \textbf{Detector maturity} \\
\midrule
Business truth & 125 & False-green business logic & HLT-008 / fast & Domain invariants, property tests, golden workflows, DB constraints, incident replay tests. & evidence-routed \\
Business truth & 100 & Authorization or data-isolation drift & HLT-022 / db & Role matrix tests, negative authz tests, owner-map routing, RLS where useful. & evidence-routed \\
Business truth & 80 & Idempotency or side-effect duplication & HLT-008 / release & Idempotency keys, replay tests, outbox/workflow proofs, double-charge negative cases. & evidence-routed \\
Security & 80 & Plausible insecure generated code & HLT-009 / security & SAST/SCA, secure-code rules, threat-model checklist, negative tests. & heuristic \\
Security & 80 & Secrets and identity sprawl & HLT-010 / security & Push protection, local/CI secret scans, redacted logs, credential inventory, transcript/artifact scan. & deterministic \\
Security & 75 & Customer-data or PII leakage & HLT-010 / security & Data classification, prompt/transcript redaction, retention policy, vector-store review, privacy tests. & governance-backed \\
Security & 75 & Excessive agent agency & HLT-012 / security & Least-privilege profiles, destructive-command gates, workspace-only writes, permission receipts. & heuristic \\
Security & 75 & Prompt injection and hostile context & HLT-011 / security & Source-trust hierarchy, untrusted-content isolation, tool-call validation, no secrets in context. & heuristic \\
Security & 60 & Improper AI/tool output handling & HLT-011 / security & Schema validation, output encoding, command sandboxing, no direct execution of generated text. & evidence-routed \\
Security & 45 & Dependency and supply-chain drift & HLT-016 / security & Dependency rationale, lockfile diff review, OSV/SCA, SBOM, provenance checks. & evidence-routed \\
Verification & 64 & False-green tests & HLT-008 / fast & Reviewer-owned acceptance criteria, mutation/fault checks, semantic assertions, red/green evidence. & evidence-routed \\
Context/setup & 64 & Context retrieval failure & HLT-015 / fast & Short root router, owner/test maps, local rules, symbol search, command-output filtering. & research \\
Contracts/data & 60 & Contract and DTO drift & HLT-007 / contract & Generated clients only, contract drift lane, generated-zone lock, compatibility tests. & heuristic \\
Contracts/data & 60 & Wrong-layer persistence & HLT-006 / db & DB access policy, adapters-only SQL, migration and constraint tests, forbidden imports. & heuristic \\
Contracts/data & 60 & Destructive migration or data-loss hazard & HLT-021 / db & Rollback/down plan, staged backfill, lock timeout, constraint proof, production rehearsal. & deterministic \\
Context/setup & 48 & Environment/setup hallucination & HLT-015 / fast & One-command setup, pinned toolchain, devcontainer where useful, setup proof lane. & evidence-routed \\
Verification & 48 & Pixel/UI regression & HLT-013 / web & Storybook states, Playwright screenshots, geometry rules, baseline governance, trace artifacts. & evidence-routed \\
Verification & 48 & Accessibility regression & HLT-014 / web & axe/ARIA checks, keyboard journeys, focus/target geometry, expert review for semantic gaps. & evidence-routed \\
Verification & 48 & Model, prompt, or eval drift & HLT-008 / ai-eval & Versioned prompts/models, eval baselines, golden datasets, provider-change receipts. & research \\
Verification & 48 & Release identity or artifact drift & HLT-025 / HLT-037 / release & Version source, changelog, immutable tag, CI-built artifact, checksum/SBOM/provenance, rollback evidence. & deterministic \\
Repair & 45 & Observability and repair opacity & HLT-017 / observability & Problem details, stable error codes, OpenTelemetry attributes, correlation IDs, repair hints. & evidence-routed \\
Entropy & 45 & Dead-language normalization & HLT-001 / fast & Banned marker scan, scoped allowlist, owner, expiry, waiver record. & deterministic \\
Entropy & 40 & Orphan or dead reachable code & HLT-001 / fast & Reachability checks, owner review, deletion receipts, tests proving replacement path. & research \\
Entropy & 36 & Performance, memory, or concurrency regression & HLT-018 / release & Benchmarks, query-plan checks, load smoke tests, allocation budgets, tracing. & evidence-routed \\
Entropy & 36 & Mega-file or mega-function sprawl & HLT-008 / fast & LOC/function caps, split by owner/use case, focused tests before behavior edits. & heuristic \\
Contracts/data & 30 & Generated-zone mutation & HLT-002 / contract & Generated-zone manifest, source contract change, regeneration command, drift check. & deterministic \\
Context/setup & 30 & Documentation or instruction drift & HLT-015 / audit & Canonical manifest, generated mirrors, instruction lint, short root guidance. & heuristic \\
\bottomrule
\end{tabularx}
\end{table*}

\begin{figure*}[t]
\centering
\begin{minipage}{0.36\textwidth}
\centering
\begin{tikzpicture}[scale=1.18]
\def\r{1.55}
\fill[JKBad] (0,0) -- (90:\r) arc (90:-25:\r) -- cycle;
\fill[orange!70] (0,0) -- (-25:\r) arc (-25:-97:\r) -- cycle;
\fill[teal!62] (0,0) -- (-97:\r) arc (-97:-144:\r) -- cycle;
\fill[JKBlue!68] (0,0) -- (-144:\r) arc (-144:-191:\r) -- cycle;
\fill[yellow!70!black] (0,0) -- (-191:\r) arc (-191:-227:\r) -- cycle;
\fill[JKNeutral!85!black] (0,0) -- (-227:\r) arc (-227:-259:\r) -- cycle;
\fill[purple!55] (0,0) -- (-259:\r) arc (-259:-270:\r) -- cycle;
\foreach \angle/\label in {32/1,-61/2,-121/3,-168/4,-209/5,-243/6,-264/7} {
  \node[font=\scriptsize\bfseries] at (\angle:1.02) {\label};
}
\foreach \angle in {90,-25,-97,-144,-191,-227,-259,-270} {
  \draw[white,line width=0.7pt] (0,0) -- (\angle:\r);
}
\draw[black!45,line width=0.4pt] (0,0) circle (\r);
\node[align=center] at (0,0) {\scriptsize Jankurai\\\scriptsize RPN share};
\end{tikzpicture}
\end{minipage}
\hfill
\begin{minipage}{0.60\textwidth}
\scriptsize
\begin{tabularx}{\linewidth}{L{0.05\linewidth} L{0.09\linewidth} L{0.20\linewidth} Y}
\toprule
\JKTableHeader
\textbf{\#} & \textbf{Share} & \textbf{TLR bucket} & \textbf{Included faults} \\
\midrule
1 & 31\% & Security & Generated-code flaws, secrets, PII, prompt injection, agency, tool output, supply chain. \\
2 & 19\% & Business truth & False-green business logic, authorization/data isolation, idempotency. \\
3 & 13\% & Contracts/data & Contract drift, wrong-layer persistence, destructive migrations, generated mutation. \\
4 & 16\% & Verification & False-green tests, pixel/UI, accessibility, release identity, prompt/model/eval drift. \\
5 & 10\% & Entropy & Dead language, orphan code, mega files/functions, performance/concurrency. \\
6 & 9\% & Context/setup & Context retrieval, setup hallucination, instruction drift. \\
7 & 3\% & Repair & Observability and repair opacity. \\
\bottomrule
\end{tabularx}
\end{minipage}
\caption{TLR means Top-Level Risk. Shares are computed from Table~\ref{tab:fault-taxonomy}: each row contributes its RPN to one TLR bucket, then bucket totals are rounded to whole percentages. The chart is a Jankurai 0.8.0 policy-priority model, not an incident-frequency measurement.}
\label{fig:tlr-pie}
\end{figure*}

\begin{table*}[t]
\caption{Common non-defensible toolchain behaviors and the evidence Jankurai expects before merge.}
\label{tab:toolchain-bad-behavior}
\centering
\scriptsize
\begin{tabularx}{\textwidth}{L{0.13\textwidth} Y L{0.18\textwidth} Y}
\toprule
\JKTableHeader
\textbf{Technology} & \textbf{Non-defensible behavior} & \textbf{Rule family} & \textbf{Required proof} \\
\midrule
Rust & Hiding memory, aliasing, thread-safety, or shell-execution assumptions behind \texttt{unsafe}, unchecked APIs, broad lint suppression, or dynamic command text. & \texttt{HLT-029}; security/fast & Local \texttt{SAFETY:} proof, narrow API contract, focused negative tests, and security lane when execution boundaries move. \\
SQL/PostgreSQL & Building executable SQL strings, shipping destructive migrations without rollback/backup evidence, or relying on app-only data invariants. & \texttt{HLT-030}; \texttt{HLT-021}; db & Parameter binding, constraints, migration safety receipt, row-count/lock plan, rollback or forward-repair path, and DB proof. \\
TypeScript & Using \texttt{any}, double casts, disabled strictness, unchecked JSON, raw DOM HTML, shell, or SQL shortcuts at trust boundaries. & \texttt{HLT-031}; \texttt{HLT-023}; fast/security & Runtime decoder or schema, strict compiler lane, boundary tests for bad input, and generated contract evidence. \\
Python & Executing dynamic code, unsafe deserialization, shelling out with strings, owning product truth, or touching production DB paths without an approved exception. & \texttt{HLT-033}; \texttt{HLT-005}; contract/security & Dated exception, typed service boundary, safe loaders, argv-only execution, parameterized DB access, and containment proof. \\
Docker & Running privileged or host-namespace containers, baking secrets into layers, using mutable images, or installing remote scripts without verification. & \texttt{HLT-032}; security & Least-privilege container profile, pinned digests, checksum/signature evidence, secret-free build context, and image/security scan receipt. \\
CI & Running untrusted code with privileged tokens, broad workflow permissions, mutable actions, secret-printing debug output, or nonblocking security scans. & \texttt{HLT-034}; \texttt{HLT-020}; security/ci & Pinned actions, minimum permissions, trusted/untrusted job separation, blocking security jobs, artifact provenance, and workflow diff review. \\
Git & Automating hard resets, broad staging, stash-based hidden state, force pushes, destructive ref edits, or verification bypass. & \texttt{HLT-035}; audit & Explicit path list, clean status receipt, no hidden state, reviewable ref operation, and exact staged-path evidence. \\
Git tools/hooks/agents & Letting hooks, MCP tools, skills, or coding agents mutate broad repository state without provenance, scope, or receipt. & \texttt{HLT-024}; \texttt{HLT-035}; audit/security & Tool identity, least-privilege permission profile, supply-chain review, changed-path receipt, and rollback/stop conditions. \\
Release & Mutating release tags or assets, skipping verification, publishing mutable \texttt{latest}-only artifacts, packaging secrets, or publishing without integrity evidence. & \texttt{HLT-025}; \texttt{HLT-037}; release & Version source, changelog, protected tag or commit identity, CI release receipt, checksum/SBOM/provenance/signature evidence, rollback plan. \\
\bottomrule
\end{tabularx}
\end{table*}

\begin{center}
\fbox{\begin{minipage}{0.93\columnwidth}
\textbf{Highest-priority controls.} A Jankurai-conformant repo should block merges when a changed path has no proof lane, a public boundary has handwritten contract mirrors, durable truth moves outside the declared durable-truth owner, database constraint layer, or approved domain/application boundary, security or secret scans are absent for risky changes, agent permissions are broader than the lane requires, user-facing UI lacks rendered proof, or \texttt{fallback}/\texttt{todo}/\texttt{stub} ships without a scoped waiver.
\end{minipage}}
\end{center}

Known solutions are stronger for some rows than others. Secret scanning, SCA, contract generation, DB constraints, Playwright screenshots, axe checks, and generated-zone manifests are known controls. Model/prompt drift, AI visual triage, context retrieval failure, and human-label learning are still emerging controls; Jankurai treats them as soft findings until teams can attach deterministic evidence. The taxonomy turns insult into audit surface. A finding is useful only when it names the rule, points at exact evidence, gives a bounded repair, and can be serialized into an agent queue.
